\chapter{Diretrizes Metodológicas do Ensino}

O curso de \ppccurso da Universidade Federal de Uberlândia adota uma concepção pedagógica centrada na formação integral do(a) estudante, combinando o desenvolvimento de conhecimentos técnico-científicos com habilidades práticas, atitudes éticas e capacidades reflexivas. As diretrizes metodológicas são pautadas por uma abordagem orientada por competências (\textit{Competency-Based Learning}), que valoriza o protagonismo discente, a interdisciplinaridade e a articulação entre teoria e prática.

\section{Princípios Metodológicos}

As práticas de ensino-aprendizagem são orientadas pelos seguintes princípios:

\begin{itemize}
    \item Integração entre teoria e prática, por meio de projetos, estudos de caso, práticas laboratoriais, desafios de engenharia e atividades extensionistas;
    \item Interdisciplinaridade, promovendo conexões entre áreas da ciência da computação, engenharia elétrica, eletrônica, controle e automação, ciências humanas e sociais aplicadas;
    \item Flexibilidade curricular, com espaço para percursos formativos personalizados, disciplinas optativas e reconhecimento de saberes adquiridos em experiências acadêmicas e profissionais;
    \item Aprendizagem ativa, com incentivo ao uso de metodologias como sala de aula invertida\footnote{\textit{Flipped Classroom}, modelo pedagógico em que a lição de casa tradicional se torna o aprendizado em casa, e as atividades em sala de aula são focadas em exercícios e discussões mais aprofundadas.}, aprendizagem baseada em projetos (PBL), resolução de problemas (ABP), estudos dirigidos e simulações computacionais;
    \item Aprendizagem ativa, com incentivo ao uso de metodologias como sala de aula invertida, aprendizagem baseada em projetos (PBL) e resolução de problemas (ABP), estudos dirigidos e simulações computacionais -- tais abordagens, conforme destacado pelo relatório \textit{Engineering for Sustainable Development}~\cite{unesco_eng4susdev_2021}, são fundamentais para que os estudantes desenvolvam competências para analisar e solucionar problemas de alta complexidade, integrando saberes de forma interdisciplinar e colaborativa em prol de um desenvolvimento sustentável;
    \item Tecnologias digitais na educação, aproveitando plataformas virtuais, ambientes de codificação online, ferramentas colaborativas e sistemas de avaliação digital;
    \item Educação híbrida, integrando atividades presenciais e remotas de forma planejada, conforme a natureza dos componentes curriculares e as necessidades dos(as) estudantes;
    \item Formação continuada, com incentivo à aprendizagem ao longo da vida, à autonomia intelectual e à curiosidade investigativa.
\end{itemize}

\section{Uso de Educação a Distância (EaD)}
\label{sec:ead}

O curso de \ppccurso da UFU adota a Educação a Distância (EaD) como uma ferramenta complementar de ensino, com o objetivo de promover metodologias ativas, acesso ampliado a conteúdos, flexibilidade pedagógica e inovação tecnológica. Essa prática está em plena conformidade com o Decreto nº 12.456, de 19 de maio de 2025~\cite{mec_decreto_ead_12456_2025}, e com a Resolução CNE/CES nº 5/2016~\cite{cne_ces_5_2016}.

Em cursos de graduação presenciais, como é o caso deste, é permitida a inclusão de até \textbf{30\% da carga horária total} do curso com atividades síncronas ou assíncronas realizadas por EaD, conforme o Art. 10 do referido Decreto.

\subsection{Definições e diretrizes fundamentais}

As modalidades de EaD previstas neste PPC incluem:

\begin{itemize}
    \item Atividades síncronas: realizadas com a presença simultânea de docentes e discentes por meio de recursos de áudio e vídeo (ex: videochamadas, salas virtuais ao vivo);
    \item Atividades assíncronas: realizadas em tempos distintos por docentes e discentes (ex: fóruns, videoaulas, quizzes, tarefas em AVA, sala de aula invertida);
    \item Atividades síncronas mediadas: síncronas com grupo de até 70 estudantes por docente ou mediador, com controle de frequência obrigatório;
    \item Atividades presenciais obrigatórias: que podem ocorrer na sede da UFU, em laboratórios, em ambientes profissionais ou espaços de extensão, conforme o plano pedagógico.
\end{itemize}

\subsection{Planejamento e supervisão}
\label{sec:planejamento_ead}

Toda unidade curricular com carga EaD será planejada e supervisionada por docente responsável (professor regente), podendo haver apoio de mediadores pedagógicos e tutores, conforme previsão legal. Os materiais didáticos seguirão os princípios de:
\begin{itemize}
    \item alinhamento com os objetivos de aprendizagem e as competências da disciplina;
    \item acessibilidade e diversidade de recursos;
    \item compatibilidade com as Diretrizes Curriculares Nacionais e os referenciais da UFU.
\end{itemize}

As atividades EaD do curso serão desenvolvidas preferencialmente nas plataformas institucionais Moodle UFU ou Microsoft Teams, que oferecem suporte ao acompanhamento das atividades síncronas e assíncronas, registro acadêmico, interação pedagógica e gestão de conteúdos digitais.

\textbf{Disciplinas originalmente presenciais poderão adotar}, de forma planejada e registrada em plano de ensino, \textbf{uma fração de sua carga horária em formato EaD}, caracterizando-se como componentes semipresenciais. Essa fração, \textbf{limitada a até 20\% da carga horária total da disciplina}\footnote{Respeitando o limite de 30\% da carga horária total do curso para atividades em EaD, conforme o Decreto nº 12.456/2025.}, poderá ser distribuída ao longo do semestre letivo, desde que respeite os princípios pedagógicos e os critérios de avaliação e frequência definidos no plano de ensino. As atividades EaD inseridas nesses componentes poderão incluir semanas assíncronas com tarefas digitais, videoaulas, fóruns e leituras, bem como encontros síncronos mediados por plataformas virtuais.

\subsection{Avaliação das atividades EaD}

As avaliações de unidades curriculares com EaD seguirão os princípios gerais do curso, devendo observar:

\begin{itemize}
    \item obrigatoriedade de avaliações presenciais para atividades EaD que representem fração significativa da carga horária;
    \item desenvolvimento de habilidades de análise e síntese discursiva, compondo ao menos 1/3 da nota, salvo exceções práticas (Art. 23, §1º, III)~\cite{mec_decreto_ead_12456_2025};
    \item garantia da autenticidade e identificação dos(as) estudantes nas avaliações síncronas e presenciais.
\end{itemize}

\subsection{Compromisso institucional}

A UFU, como Instituição de Educação Superior pública federal, está automaticamente credenciada para a oferta complementar de atividades EaD, conforme Art. 15 do Decreto nº 12.456/2025. O curso de \ppccurso reafirma seu compromisso com:

\begin{itemize}
    \item o uso responsável, pedagógico e inovador das tecnologias digitais;
    \item o respeito à legislação vigente, aos limites estabelecidos e à identidade do curso;
    \item a valorização da interação humana, do engajamento ativo e do aprendizado de excelência.
\end{itemize}

\section{Integração com Pesquisa, Extensão e Vivência Profissional}

A metodologia do curso articula ensino, pesquisa e extensão como dimensões indissociáveis da formação universitária. Os(as) estudantes são incentivados(as) a participar de projetos de iniciação científica, programas de extensão, atividades em laboratórios especializados, empresas juniores, PET, hackathons, maratonas de programação, estágios e cooperação internacional.

Além disso, o curso valoriza as aprendizagens desenvolvidas em contextos extracurriculares, por meio de mecanismos formais de aproveitamento de estudos e competências adquiridas fora da sala de aula, conforme normativas internas da UFU.

\section{Avaliação do Processo de Ensino-Aprendizagem}

A avaliação no curso é formativa, diagnóstica e contínua, buscando não apenas aferir resultados, mas orientar o desenvolvimento progressivo das competências previstas. As estratégias de avaliação são variadas e contextualizadas, incluindo:

\begin{enumerate}
    \item provas escritas e orais;
    \item projetos e trabalhos em grupo;
    \item atividades práticas e laboratoriais;
    \item autoavaliações e coavaliações;
    \item apresentações técnicas e relatórios;
    \item portfólios e produções técnicas autorais.
\end{enumerate}

A coerência entre métodos de ensino, objetivos da disciplina, competências visadas e formas de avaliação é elemento fundamental da qualidade do processo formativo.

\section{Uso Ético e Consciente de IA no Ensino-Aprendizagem}
\label{sec:IA-UFU}

O curso de \ppccurso da UFU reconhece a Inteligência Artificial (IA), em especial a IA Generativa, como uma tecnologia transformadora com profundo impacto no ensino, na pesquisa e na prática profissional da engenharia. Alinhado ao Guia de Princípios para o Uso Ético e Responsável de IA da UFU~\cite{ufu_guia_ia_2025} e às discussões contemporâneas sobre o tema, este PPC não busca restringir o uso dessas ferramentas, mas sim promover uma cultura de letramento em IA, capacitando os(as) discentes a se tornarem curadores conscientes da cognição, em vez de meros consumidores de tecnologia.

Nossa abordagem se fundamenta na premissa de que o verdadeiro risco da IA não é a substituição do intelecto humano, mas a abdicação voluntária de nossa capacidade de pensar criticamente. Portanto, as seguintes diretrizes orientam o uso de IA no percurso formativo:

\subsection{Princípios Éticos e Integridade Acadêmica}

\begin{enumerate}
    \item \textbf{Supervisão Humana e Responsabilidade Final:} Toda e qualquer utilização de IA deve ser rigorosamente supervisionada. O(a) discente é o(a) autor(a) e responsável final por qualquer conteúdo que apresente, cabendo-lhe a verificação da acurácia, a validação das fontes e a originalidade do trabalho. A IA é uma ferramenta de apoio, não um substituto para o discernimento humano.

    \item \textbf{A Distinção entre Plausibilidade e Verdade:} Os modelos de IA Generativa são projetados para sintetizar conteúdo plausível, que se assemelha a padrões aprendidos, mas não necessariamente verdadeiro ou factualmente correto. É imperativo que os(as) estudantes compreendam que as informações geradas por IA podem conter ``alucinações algorítmicas'' e vieses. A checagem de fatos e a busca por fontes primárias confiáveis são competências essenciais e não negociáveis.

    \item \textbf{Transparência e Autoria:} Em conformidade com o princípio da integridade acadêmica, é obrigatório explicitar, de forma clara e transparente, toda e qualquer utilização de ferramentas de IA no processo de concepção, desenvolvimento ou redação de trabalhos acadêmicos. A forma de citação e declaração de uso será orientada pelos docentes em cada componente curricular.

    \item \textbf{Privacidade e Segurança:} É vedada a inserção de dados pessoais sensíveis, informações estratégicas da universidade ou dados de terceiros em plataformas de IA públicas que não garantam os requisitos adequados de segurança e proteção jurídica, em observância à Lei Geral de Proteção de Dados (LGPD).
\end{enumerate}

\subsection{O Papel da IA como Ferramenta de Aprendizagem}

Este curso incentiva o uso consciente da IA como uma poderosa aliada para potencializar a aprendizagem e a criatividade. Encorajamos os(as) estudantes a enxergarem a IA não como um oráculo, mas como um colega de equipe incansável, que pode assumir múltiplos papéis no processo de estudo.

\begin{enumerate}
    \item \textbf{A IA como ``Monitor de Disciplinas em Tempo Integral'':} Os(as) discentes podem utilizar a IA como um assistente de estudos personalizado, disponível a qualquer momento para:
    \begin{itemize}
        \item Explicar conceitos complexos com diferentes analogias.

        \item Sugerir roteiros de estudo e materiais complementares.

        \item Gerar problemas e exercícios para praticar e fixar o conhecimento.

        \item Auxiliar na depuração de códigos e na compreensão de erros.
        Contudo, reforça-se que este ``monitor'' é um ponto de partida. A validação final do conhecimento e o aprofundamento conceitual devem ser sempre realizados com os docentes e os materiais de referência da disciplina.
    \end{itemize}

    \item \textbf{Combatendo o ``Viés da Primeira Ideia'':} A criatividade, em sua essência, é ``fazer mais do que a primeira coisa que se pensa''~\cite{utley_2025}. A IA é uma ferramenta excepcional para superar a tendência humana de se fixar na primeira solução aparente. Incentiva-se o uso da IA para:
    \begin{itemize}
        \item Gerar um grande volume de ideias e abordagens alternativas para um problema (\textit{brainstorming}).

        \item Explorar diferentes perspectivas e refutar hipóteses iniciais.

        \item Ampliar o repertório criativo, forçando o(a) estudante a ir além do óbvio.
    \end{itemize}

    \item \textbf{A Importância do Contexto e da Curadoria:} A qualidade da interação com a IA é diretamente proporcional à qualidade do contexto fornecido pelo usuário. Para obter resultados valiosos, o(a) estudante deve atuar como um(a) curador(a) da cognição, fornecendo \textit{briefings} detalhados, questionando as respostas, refinando os comandos (\textit{prompts}) e, principalmente, adicionando sua visão crítica e seu conhecimento de mundo ao produto final. A IA amplifica a capacidade do usuário; ela não a cria do zero~\cite{utley_2025}.
\end{enumerate}

\subsection{Formação e Letramento em IA}

O curso se compromete a promover o letramento em IA de forma contínua e transversal. A disciplina de Introdução à Engenharia de Computação abordará formalmente estas diretrizes, preparando os ingressantes para uma jornada acadêmica consciente. Adicionalmente, os docentes irão, em seus respectivos componentes curriculares, orientar as boas práticas e os limites do uso de IA, fomentando uma postura de aprendizado e adaptação contínuos, essencial para a atuação profissional nesta nova era.